\documentclass{ximera}

\title{Linearization Activity Solution Guide}
\author{MATH 425: Calculus I}

\begin{document}
\begin{abstract}
   Working with the peers in your group, solve the following problems. Make sure to show and justify all your work. Make sure everyone in the group understands the solution and participates. Be prepared to report your answers to the whole class. 
\end{abstract}
\maketitle


\begin{exercise}
   Let $c$ be a nonzero number. Consider the function $f(x)=\sin(cx)$ and the point $(0,0)$. \begin{enumerate}
      \item Find the linearization of $f$ at $0$.\\
        \textcolor{blue}{ $f(x)\approx c(x-0)+0=cx$}
        
        \item Find an estimate for $f(-0.1)$. Find an estimate for $f(0.123)$.\\
        \textcolor{blue}{ $f(-0.1)\approx c(-0.1)=-0.1c$\\
        $f(0.123)\approx 0.123c$}
        
        \item How might this linearization help you find the limit $\lim_{x\to 0}\frac{\sin(cx)}{cx}$?\\
        \textcolor{blue}{From (b), it seems as though $\sin(cx)\approx cx$ when $x$ is close to $0$. So, $\lim_{x\to 0}\frac{\sin(cx)}{cx}\approx \lim_{x\to 0}\frac{cx}{cx}=1$.}
   \end{enumerate}
\end{exercise}
 
\begin{exercise}
   Consider the functions $g(x)=2x^{2}-7x$ and $h(x)=x^{3}-4x^{2}+2x$. 
   \begin{enumerate}
       \item Show that these functions have the same linearization $L(x)$ at $a=3$.
        \textcolor{blue}{ $g(3)=2(3^2)-7(3)=18-21=-3$\\
        $g'(3)=4(3)-7=5$\\
        $g(x)\approx 5(x-3)-3$\\
        $h(3)=3^3-4(3^2)+2(3)=27-36+6=-3$\\
        $h'(3)=3(3^2)-8(3)+2=27-24+2=5$\\
        $h(x)\approx 5(x-3)-3$\\
        So $L(x)=5(x-3)-3$}
    
        
        \item Compute $L(2.8)$, $g(2.8)$, and $h(2.8)$. Which does $L(2.8)$ better approximate: $g(2.8)$ or $h(2.8)$? \\
        \textcolor{blue}{ $L(2.8)=5(2.8-3)-3=5(-.2)-3=-4$\\
        $g(2.8)=-3.92$\\
        $h(2.8)=-3.808$\\
        $g(x)$ is better approximated by $L(x)$.}
        
        \item Think about your responses for (b). Why do think this occurs? Graph the functions $L(x)$, $g(x)$, and $h(x)$. What do you notice?
        \textcolor{blue}{ $h(x)$ stays closer to the tangent line $L(x)$, as it is less ``turn-y''.  It seems like the fewer turns (at least nearby $a$), the better the approximation is.}
   \end{enumerate}
\end{exercise}

\begin{exercise}
   Suppose that the area of a circular oil spill is expanding at a rate of $15\pi \frac{ft^2}{min}$.  How fast is the radius of the spill changing at the moment the radius is 6 ft? (Let $r(t)$ be the radius at time $t$.) \\
   {\textcolor{blue} $A(t)=\pi r(t)^2$ }\\
   $A'(t)=2\pi r(t)r'(t)$\\
   We know $A'(t)=15$ and $r(t)=6$$\\
   $15=2\pi (6)r'(t)$\\
   So $r'(t)=\frac{15}{12\pi}$ 
\\
One way to check your answer  to \#4 is to use a linear approximation to estimate the radius a short time after it is 6 ft. In the previous problem, you should have found $\frac{dr}{dt}$ when the radius is $6$ ft. Let's call that time $t=0$, so that you found $r'(0).$ And $r(0)=6$ ft. Use a linear approximation of $r(t)$ at $t=0$ to estimate the radius at $t=0.2$ min. The exact value of the radius at this time is $\sqrt{39}\approx 6.245$ ft.\\
{\textcolor{blue} Linearization at time $t$: $L(x)=r'(0)(x-0)+r(0)$}\\
$L(0.2)=\frac{15}{12\pi}(0.2)+r(\sqrt{39})\approx r(0.2)$
\end{exercise}


\end{document}
